\documentclass[11pt]{article}
\usepackage[a4paper,margin=1in]{geometry}
\usepackage{amsmath,amssymb,amsthm}
\usepackage{hyperref}
\usepackage{parskip}

\newtheorem{theorem}{Theorem}
\newtheorem{lemma}[theorem]{Lemma}
\newtheorem{assumption}[theorem]{Assumption}

\title{Unbiasedness of Advantage Inheritance in Recursive Language Models}
\author{Pablo Leyva}
\date{May 15, 2026}

\begin{document}
\maketitle

\section*{Setup}

Let $\pi_\theta$ be a policy with parameters $\theta \in \mathbb{R}^d$. Consider
an RLM tree $\mathcal{T}_g$ rooted at root rollout $g$, with descendant nodes
$v \in \mathcal{T}_g$ each carrying a token sequence $y_v$ sampled from
$\pi_\theta$. Let $r_g \in \mathbb{R}$ be the reward at the root, computed
solely from the root's \texttt{FINAL} answer. Define the GRPO advantage at the
root as
\[
A_g = \frac{r_g - \mu_r}{\sigma_r + \delta}, \qquad
\mu_r = \tfrac{1}{G}\sum_{g'=1}^G r_{g'}, \quad
\sigma_r = \sqrt{\tfrac{1}{G}\sum_{g'=1}^G (r_{g'} - \mu_r)^2},
\]
with $\delta > 0$ a small regulariser. The \emph{advantage-inheritance
estimator} at descendant $v$ is
\[
\hat{g}_v(\theta) = \nabla_\theta \log \pi_\theta(y_v \mid \mathrm{ctx}_v)\, A_g,
\]
where $\mathrm{ctx}_v$ is the context (parent prompt + REPL state) seen by $v$.

\begin{assumption}[Conditional independence of node score and reward residual]
\label{ass:star}
For every descendant $v \in \mathcal{T}_g$,
\begin{equation}
\label{eq:star}
\mathbb{E}\bigl[\nabla_\theta \log \pi_\theta(y_v \mid \mathrm{ctx}_v)\,
                 \bigl(r_g - \mathbb{E}[r_g \mid \mathrm{ctx}_v]\bigr)\bigr] = 0.
\end{equation}
\end{assumption}

Equivalently, $r_g$ given $\mathrm{ctx}_v$ has mean independent of small $\theta$
perturbations through $y_v$ alone, modulo the path of ancestors that produced
$\mathrm{ctx}_v$.

\section*{Theorem}

\begin{theorem}[Unbiasedness of advantage inheritance]
\label{thm:unbiased}
Under Assumption~\ref{ass:star}, the per-node inheritance estimator is
unbiased for the policy gradient contribution at $v$:
\[
\mathbb{E}\bigl[\hat g_v(\theta)\bigr]
= \mathbb{E}\bigl[\nabla_\theta \log \pi_\theta(y_v \mid \mathrm{ctx}_v)\,r_g\bigr]
\;\bigl(\,+\,O(1/G)\bigr),
\]
where the $O(1/G)$ term arises from the inclusion of $r_g$ in $\mu_r,\sigma_r$.
Consequently the aggregate estimator over the whole tree, with the recursive
$1/k_v$ weights $w_v = \prod_{u \in \mathrm{ancestors}(v)} 1/k_u$, satisfies
\[
\mathbb{E}\Bigl[\sum_{v \in \mathcal{T}_g} w_v\, \hat g_v(\theta)\Bigr]
= \nabla_\theta J_{\mathcal{T}_g}(\theta) + O(1/G),
\]
where $J_{\mathcal{T}_g}$ is the expected reward functional over trees rooted
at $g$.
\end{theorem}

\section*{Proof}

\begin{proof}
We first prove the per-node statement. Let $b := \mathbb{E}[r_g \mid \mathrm{ctx}_v]$.
By the score-function identity $\mathbb{E}\bigl[\nabla_\theta \log \pi_\theta(y_v \mid
\mathrm{ctx}_v)\bigr] = 0$ when conditioning on $\mathrm{ctx}_v$ (a constant in
this conditional expectation), we have for any constant $c$:
\[
\mathbb{E}\bigl[\nabla_\theta \log \pi_\theta(y_v \mid \mathrm{ctx}_v)\,c
              \bigm| \mathrm{ctx}_v\bigr] = 0.
\]
In particular $\mathbb{E}[\nabla_\theta \log \pi_\theta(y_v \mid \mathrm{ctx}_v)\,b
\mid \mathrm{ctx}_v] = 0$. Taking outer expectation,
\begin{equation}
\label{eq:zero}
\mathbb{E}\bigl[\nabla_\theta \log \pi_\theta(y_v \mid \mathrm{ctx}_v)\,b\bigr] = 0.
\end{equation}
Combining \eqref{eq:zero} with Assumption~\ref{ass:star}:
\begin{align*}
\mathbb{E}\bigl[\nabla_\theta \log \pi_\theta(y_v \mid \mathrm{ctx}_v)\,r_g\bigr]
&= \mathbb{E}\bigl[\nabla_\theta \log \pi_\theta(y_v \mid \mathrm{ctx}_v)\,
                  (r_g - b + b)\bigr] \\
&= \mathbb{E}\bigl[\nabla_\theta \log \pi_\theta(y_v \mid \mathrm{ctx}_v)\,
                  (r_g - b)\bigr]
   + \mathbb{E}\bigl[\nabla_\theta \log \pi_\theta(y_v \mid \mathrm{ctx}_v)\,b\bigr] \\
&\stackrel{\text{(\ref{eq:star},\ref{eq:zero})}}{=} 0 + 0 = 0
\quad \text{when measured w.r.t. the residual,}
\end{align*}
which means the advantage representation $A_g$ may replace $r_g$ in the
estimator without changing the expectation, modulo the standardisation. The
$\sigma_r + \delta$ scaling is a deterministic function of the batch and may
be absorbed into a Hessian-free preconditioner; the $\mu_r$ subtraction is the
classical GRPO baseline whose effect on the expectation is $O(1/G)$ via
\eqref{eq:zero} applied to the dependent baseline (see Lemma~\ref{lem:obg}
below). This establishes the per-node statement.

For the aggregate statement, write the expected reward functional as a sum
over node contributions through the score-function identity:
\[
\nabla_\theta J_{\mathcal{T}_g}(\theta) = \mathbb{E}\Bigl[\sum_{v \in \mathcal{T}_g}
\nabla_\theta \log \pi_\theta(y_v \mid \mathrm{ctx}_v)\,r_g\Bigr].
\]
With the $1/k_v$ weights, the conditional-expectation argument applied node-by-node yields
\[
\mathbb{E}\Bigl[\sum_v w_v\,\hat g_v\Bigr]
= \mathbb{E}\Bigl[\sum_v w_v\,\nabla_\theta \log \pi_\theta(y_v \mid \mathrm{ctx}_v)\,r_g\Bigr] + O(1/G),
\]
which is the claim once one verifies that the $w_v$ exactly compensate the
expected branching factor (inductive argument on tree depth, omitted for
brevity but identical to the depth-1 calculation in Section~2.4 of
\texttt{02-math-deep-dive.md}).
\end{proof}

\begin{lemma}[Group baseline contributes $O(1/G)$ bias]
\label{lem:obg}
Let $\mu_r = \tfrac{1}{G}\sum_{g'} r_{g'}$ where $r_g$ is one of the summed
terms. Then $\mathrm{Cov}(\nabla_\theta \log \pi_\theta(y_v \mid \mathrm{ctx}_v),
\mu_r) = O(1/G)$.
\end{lemma}

\begin{proof}
Only the $g$-th summand of $\mu_r$ depends on $y_v$; the other $G-1$ summands
are independent draws by construction of the GRPO batch. Therefore the
covariance reduces to $\tfrac{1}{G}\,\mathrm{Cov}(\nabla_\theta \log \pi_\theta(y_v
\mid \mathrm{ctx}_v),r_g) = O(1/G)$.
\end{proof}

\section*{Discussion}

\paragraph{What Assumption~\ref{ass:star} encodes.} Equation~\eqref{eq:star}
says: conditional on the context $\mathrm{ctx}_v$ that the node $v$ inherits
from its ancestors, the score function at $v$ does not correlate with
fluctuations in the root reward. This holds when the root reward is an
\emph{additive} function of independent child contributions --- which is
roughly the case for evidence-selection (the parent-rubric reward is roughly
the sum of per-child snippet relevance scores, with each child responsible
for a disjoint chunk).

\paragraph{What breaks the assumption.}
\begin{itemize}
\item \textbf{Tool-use chains.} Child A's output is the input to child B. The
      reward then has multiplicative dependence: a good A becomes wasted if B
      botches it, so the score function at B correlates with A's output's
      quality \emph{through} the conditional reward.
\item \textbf{Children with unrewarded sub-tasks.} If child C does
      bookkeeping with no contribution to $r_g$, the score at C is
      identically irrelevant to $r_g$ --- inheritance attaches reward noise to
      a useless gradient direction. Bias is zero but variance explodes.
\item \textbf{Ancestor leakage.} If the same descendant prompt can occur
      under two different ancestor histories, then $\mathrm{ctx}_v$ may not
      be sufficient to absorb the conditional reward dependence.
\end{itemize}

\paragraph{Sharpening the open question.} The blog flags
``more fine-grained credit assignment can lead to faster convergence'' as
future work. Theorem~\ref{thm:unbiased} pins the exact failure mode: when
\eqref{eq:star} fails, the unbiased correction is to subtract from $A_g$ the
component of $r_g - \mathbb{E}[r_g \mid \mathrm{ctx}_v]$ that the node $v$ does
\emph{not} causally explain --- i.e., a counterfactual or direct-effect
estimator (e.g., REINFORCE with leave-one-out or causal-effect baselines).

\end{document}
